% !PRINTED 145
\weberpage{145}
Nehmen wir, was freisteht, an, dass $b \geqq a$ sei, so ist der
Ausdruck auf der rechten Seite dieser Ungleichung grösser
als $n$, und nur dann gleich $n$, wenn $b = n - 1, a = 1$ ist. In
diesem speciellen Falle kann $j = n$ werden, aber nur dann,
wenn die $n - 1$ Elemente durch $Q$ auf alle möglichen Arten
permutirt werden, was in dem Satze I. ausgesprochen ist.

Die Ausdrücke (3) für die untere Grenze von $j$ sind nichts
Anderes, als die Binomialco\"efficienten $B_a^{(n)}$, deren Bildung
sofort zeigt, dass sie bis zur Mitte hin, d. h. so lange $2a \leqq n$,
eine \gesperrt{wachsende Zahlenreihe} bilden.

Wir wollen für den weiteren Gebrauch hieraus den Schluss
ziehen:

$\alpha$) \gesperrt{Ist $a = 1$, lässt also die Gruppe $Q$ eine Ziffer
ungeändert, so ist ihr Grad ein Theiler von
$\Pi(n-1)$, ist aber $a > 1$, so ist der Grad von $Q$
$\leqq \Pi(n-2)\,\Pi(2)$.}

Der Satz, den wir ferner noch beweisen wollen, lautet nun:

\gesperrt{II. Ausser der alternirenden Gruppe giebt es keinen
transitiven und primitiven Theiler $Q$ von $P$,
dessen Index $\leqq n$ ist, ausgenommen in den
beiden Fällen $n = 4, n = 6$.}

Beim Beweise machen wir Gebrauch von den Sätzen (\S.\,153,
9., 10. des ersten Bandes), dass ein transitiver und primitiver
Theiler von $P$, der nicht die ganze alternirende Gruppe enthält,
und daher nicht mit der alternirenden oder der symmetrischen
Gruppe selbst identisch ist, keine Transposition und keine
cyklische Permutation von nur drei Ziffern enthalten kann.

Es sei also $P$ die symmetrische Permutationsgruppe der
$n$ Ziffern $0, 1, 2 \dots n - 1$ und $Q$ ein primitiver und transitiver
Theiler von $P$ vom Index $j$, der nicht die alternirende Gruppe
enthält. Es sei ferner $P$ in die Nebengruppen zerlegt:
\weq{4}{P = Q + Q\pi_1 + Q\pi_2 + \dots + Q\pi_{j-1}.}
Wir betrachten das System der Transpositionen:
\weq{5}{(0, 1), (0, 2), \dots (0, n-1),}
deren keine in $Q$ vorkommen kann. Ist nun $j < n$, so müssen
wenigstens zwei dieser Permutationen, etwa $(0,1), (0,2)$, in der-
%CONT>>
